%% ma: L12-B11-0003
%% lop: 12
%% chuong: 4
%% bai: 11
%% dang: 12.11.3
%% loai: TN4
%% muc: TH
%% dap_an: "C"
%% nguon: "(NBV)-ĐỀ SỐ 5-MỖI NGÀY 1 ĐỀ THI 2025.pdf (D05, MỖI NGÀY 1 ĐỀ - Đề số 5), Phần I Câu 9"
%% kien_thuc: "Bài 11 (nguyên hàm của hàm số lượng giác); công thức hạ bậc (lớp 11, Bài 2)"
%% trang_thai: cho-duyet
%% ly_do: "Dạng toán mới đề xuất 12.11.3: nguyên hàm của hàm lượng giác bằng công thức hạ bậc"
\begin{ex}
Nguyên hàm của hàm số $f(x)=\sin^2\dfrac{x}{2}$ là
\choice{$\dfrac{x}{2}+\dfrac{\sin x}{2}+C$}{$\dfrac{x}{2}-\dfrac{\cos x}{2}+C$}{\True $\dfrac{x}{2}-\dfrac{\sin x}{2}+C$}{$\dfrac{x}{2}+\dfrac{\cos x}{2}+C$}
\loigiai{Hạ bậc: $\sin^2\dfrac{x}{2}=\dfrac{1-\cos x}{2}$. Do đó
\[\int\sin^2\dfrac{x}{2}\,\mathrm dx=\int\dfrac{1-\cos x}{2}\,\mathrm dx=\dfrac{x}{2}-\dfrac{\sin x}{2}+C.\]}
\end{ex}
